\chapter{Cognitive Science of Concepts, Reformalized}\label{ch:cogsci}

\epigraph{The mind is not a single instrument but a workshop of instruments;
each theory of the concept is, on close inspection, a description of one of
its tools.}{}

\section*{Orientation}

The cognitive science of concepts is a notoriously fractious discipline.
Definitional, prototype, theory-theoretic, exemplar, perceptual-symbolic,
language-of-thought, embodied, conceptual-space, distributed,
tensor-binding, and now transformer-based accounts each insist on their own
ontology of what a concept \emph{is}. Each, viewed without the others, is
both insightful and incomplete. The thesis of this chapter is that the idea
algebra $(\Ideas,\compose,\unit,\rho,\sigma,\nu,\eta,\omega,\deg)$ is not a
fifteenth competitor but a \emph{translation manual}: every position
mentioned above can be reread as a coherent restriction or specialization of
the formalism, and each becomes the focus on a particular operator or
sub-algebra.

Throughout this chapter we work inside a fixed idea algebra $\IdeaAlg$
satisfying the foundational theorems of \cref{ch:foundations}. We write
$V(\Ideas)$ for the free $\RR$-vector space on the carrier $\Ideas$ and
extend $\rho:\Ideas\times\Ideas\to\RR$ bilinearly to
$\rho:V(\Ideas)\times V(\Ideas)\to\RR$ (Theorem 2). When a position
distinguishes a sub-monoid $\Ideas_C\subseteq\Ideas$ closed under
$\compose$, the restriction
$(\Ideas_C,\compose,\unit,\rho|_C,\sigma|_C,\nu|_C,\eta|_C,\omega|_C,\deg|_C)$
is again an idea algebra; its existence and equivalence type are governed
by Theorem 8 (Structural Equivalence). Each section concludes with a
\textbf{Status} clause naming the foundational theorem (or a small
corollary) that underwrites the position.

The strategy of the chapter is uniform. Each position is presented in three
movements. First, the original claim is restated faithfully, with citations
to the empirical or theoretical literature in which it was advanced. Second,
the position is translated into the operators of $\IdeaAlg$. Some positions
will pick out a sub-monoid (definitional theory, frames, theory-theory);
others will single out an operator (prototypes single out $\rho$; blending
singles out $\eta$); still others will impose a regime on the cocycle
$\omega$ (the language of thought zeroes it out, holographic reduced
representations bound it, transformers learn it). Third, the formal
status of the position is fixed by naming the foundational theorem or by
proving a small corollary in math prose. The taxonomy that emerges is not a
hierarchy of right and wrong views; it is a stratification of a common
algebraic object. The chapter closes by gluing the strata together into a
single layered architecture, \cref{thm:cogsci-layered}.

% =====================================================================
\section{The classical (definitional) theory}\label{sec:cogsci-classical}

\textbf{Original claim.} The classical theory, traceable through Aristotle,
Frege, and into mid-twentieth-century psychology, holds that a concept is a
\emph{definition}: a finite list of necessary and jointly sufficient
features. To possess the concept \textsc{bachelor} is to deploy the
intersection of the features \textsc{adult}, \textsc{male}, \textsc{human},
\textsc{unmarried}. Categorization is then crisp set membership; learning
is hypothesis testing over feature combinations \citep{Locke1690,Fodor1975}.

\formalclaim Fix a finite \emph{feature set} $F\subseteq\Ideas$ and let
$\langle F\rangle_{\compose}$ be the sub-monoid of $\Ideas$ generated by
$F\cup\{\unit\}$ under $\compose$. Imposing in addition idempotence
($f\compose f=f$) and commutativity ($f\compose g = g\compose f$) for
$f,g\in F$ presents this sub-monoid as the free \emph{Boolean meet
sub-monoid} $\mathsf{Meet}(F)\cong 2^{F}$. A classical concept is then a
homomorphism $c:\mathsf{Meet}(F)\to\{0,1\}$ given by an upward-closed
filter: an element $a\in\mathsf{Meet}(F)$ falls under $c$ iff
$\sigma(a,d_c)=d_c$ for some characteristic conjunction
$d_c\in\mathsf{Meet}(F)$.

\begin{corollary}[Crispness corollary]\label{cor:cogsci-classical}
On the Boolean sub-monoid $\mathsf{Meet}(F)$ with the canonical
$\rho_{\mathrm{cl}}(a,b)=\mathbf{1}[a\le b]$, the resonance pairing is
$\{0,1\}$-valued and the Aufhebung decomposition collapses to
$\sigma(a,b)=a\wedge b$, $\nu(a,b)=0$, $\eta(a,b)=0$. The cocycle
$\omega\equiv 0$.
\end{corollary}

\begin{proof}[Sketch]
Idempotence makes $\compose=\wedge$, hence $\eta(a,b)=0$ since no
``new content'' is generated; commutativity kills any antisymmetric
component, so $\nu=0$; uniqueness mod $\ker\rho$ from Theorem~3 then
forces $\sigma=a\wedge b$. The cocycle vanishes by Theorem~4 because the
valuation $v(a)=|a|$ on the Boolean lattice is already a homomorphism.
\end{proof}

The empirical failures of the classical theory --- typicality gradients,
borderline cases, family resemblance \citep{RoschMervis1975} --- are
precisely the failure modes of a model in which $\rho$ is forced to be
$\{0,1\}$-valued and $\eta$ vanishes.

\formalstatus The classical theory is the restriction of the formalism to
the Boolean sub-monoid generated by a finite feature set; existence and
uniqueness of the restriction is Theorem~1 (Composition) and the structural
collapse follows from \cref{cor:cogsci-classical}.

\litcomment The classical theory was not foolish; it was the correct first
guess given the prevailing model of definition inherited from logic. What
the formalism reveals is that the classical model is a single, very
restricted slice of the algebra: the slice on which composition is the
Boolean meet, on which $\rho$ takes only the values $0$ and $1$, on which
$\eta$ vanishes, and on which the cocycle $\omega$ is identically zero.
Every empirical pressure that broke the classical theory --- typicality,
graded membership, family resemblance, cross-classification, contextual
re-coordination --- arose from the parts of the algebra that this slice
discards. From this vantage the empirical defeat of the classical theory
is not a refutation of feature-based concepts but a discovery that
features are themselves graded ideas, that their composition produces
genuinely new content via $\eta$, and that the cocycle $\omega$ is not
identically zero. The classical theory survives as the limit case in
which $\rho$ is restricted to truth values, and remains the correct model
for the parts of cognition that genuinely behave that way (formal
mathematics, well-defined legal categories, certain technical
vocabularies). The formalism therefore preserves the local validity of
the classical view while exposing its global incompleteness.

% =====================================================================
\section{Rosch: prototype theory}\label{sec:cogsci-rosch}

\textbf{Original claim.} Rosch \citep{Rosch1973,RoschMervis1975} argued
that natural categories are organized around \emph{prototypes}: best
examples that compress the statistical regularities of the category. Robins
are more typical birds than penguins; chairs are more typical furniture
than televisions. Typicality is graded, reaction-time-correlated, and not
recoverable from any classical definition.

\formalclaim Let $C\subseteq\Ideas$ be a finite category with empirical
weights $w:C\to\RR_{>0}$ summing to $1$. Define the \emph{prototype} of
$C$ as the $\rho$-centroid in $V(\Ideas)$,
\[
  \pi_C \;\eqdef\; \sum_{a\in C} w(a)\,a \;\in\; V(\Ideas).
\]
The \emph{typicality} of $a\in C$ is then
$\tau_C(a)\eqdef\rho(a,\pi_C)/\sqrt{\rho(\pi_C,\pi_C)}$ when $\rho$ is
positive on the relevant subspace, and $\tau_C(a)\eqdef\rho(a,\pi_C)$
otherwise. Categorization of a fresh idea $x$ is by maximum resonance:
$\arg\max_{C}\rho(x,\pi_C)$.

\begin{corollary}[Bilinearity makes the prototype the typicality
maximizer]\label{cor:cogsci-rosch}
Within the affine hull of $C$, the function $\tau_C$ is the unique linear
functional that attains its $C$-maximum at $\pi_C$ up to scale, and is
preserved under conjugation by an invertible idea $g$:
$\tau_{gCg^{-1}}(\conj{g}{a})=\tau_C(a)$.
\end{corollary}

\begin{proof}
By Theorem~2, $\rho$ extends uniquely as a bilinear form, so
$a\mapsto\rho(a,\pi_C)$ is linear and is maximized over the simplex
$\mathrm{conv}(C)$ at the centroid (which is $\pi_C$). Theorem~6
(Conjugation) gives $\rho(\conj{g}{a},\conj{g}{b})=\rho(a,b)$, hence the
conjugacy invariance.
\end{proof}

The corollary tells us that prototype theory is not in conflict with the
formalism, but is the \emph{first-order} statistical reading of $\rho$
restricted to a category sub-monoid. Failures of prototype theory ---
asymmetric similarity, contextual flexibility \citep{Barsalou1983} --- are
the predicted residue when $\rho$ is non-symmetric (curvature $\kappa\ne 0$,
\cref{ch:foundations}).

\formalstatus Prototype theory is the bilinear-extension reading of $\rho$,
underwritten by Theorem~2 and \cref{cor:cogsci-rosch}.

\litcomment Prototype theory and the classical theory have, in the
literature, been treated as competitors. In the algebra they are
complements: the classical view fixes the sub-monoid (which features form
the category) while the prototype view fixes the operator (which
$\rho$-centroid summarizes the category). One cannot replace the other.
A category needs both a sub-monoid that says \emph{which} ideas are
candidates for membership and a centroid in $V(\Ideas)$ that says
\emph{how} graded membership ranks those candidates. The empirical
success of prototype theory at predicting reaction times and
free-listing frequencies reflects the fact that $\rho(a,\pi_C)$ is
already, in the formalism, the natural similarity-to-category measure on
$V(\Ideas)$. The formalism additionally predicts asymmetries that pure
prototype theory does not: because $\rho$ is in general non-symmetric,
$\rho(a,\pi_C)\ne\rho(\pi_C,a)$, and the difference is exactly the
curvature contribution $\kappa(a,\pi_C)$. This recovers Tversky-style
asymmetric similarity --- ``North Korea is similar to China; China is
not as similar to North Korea'' --- as a curvature phenomenon rather
than as an ad-hoc psychological observation.

% =====================================================================
\section{Murphy \& Medin: theory-theory}\label{sec:cogsci-theorytheory}

\textbf{Original claim.} \citet{MurphyMedin1985} argued that the
``coherence'' of a concept --- which features hang together, which novel
features will be projected --- cannot be explained by similarity to a
prototype or to stored exemplars. Concepts cohere because they are embedded
in tacit \emph{theories}: domain-specific causal-explanatory schemata.
\textsc{tree} hangs together because of folk biology; \textsc{money}
because of folk economics.

\formalclaim Let a \emph{domain theory} $T$ over a sub-monoid
$\Ideas_T\subseteq\Ideas$ be a triple $(\Ideas_T,\eta_T,\Phi_T)$ where
$\Ideas_T$ is closed under $\compose$, $\eta_T$ is the restriction of the
elevation operator to $\Ideas_T$, and $\Phi_T\subseteq\Ideas_T$ is a set
of \emph{theory-internal axioms} (typical causal links). The theory is
\emph{coherent} on input $a,b\in\Ideas_T$ when
\[
  \eta_T(a,b)\;\in\;\langle\Phi_T\rangle_{\compose},
\]
i.e., the synthetic content of composing $a$ with $b$ already lies in the
sub-monoid generated by $\Phi_T$. Cross-domain composition
$a\in\Ideas_T,\ b\in\Ideas_{T'}$ generically lifts $\eta(a,b)$ out of both
sub-monoids, predicting the felt incoherence of category-mistakes.

\begin{corollary}[Sub-algebraic inheritance]\label{cor:cogsci-theory}
A domain theory $T$ is itself an idea sub-algebra: the operators
$\rho_T,\sigma_T,\nu_T,\eta_T,\omega_T,\deg_T$ obtained by restriction
satisfy the axioms of \cref{ch:foundations} on $\Ideas_T$.
\end{corollary}

\begin{proof}
Closure of $\Ideas_T$ under $\compose$ ensures Theorem~1 transfers. The
remaining structure is preserved because each operator is defined
pointwise on inputs in $\Ideas_T$, and Theorem~8 (Structural Equivalence)
guarantees the restricted tuple is again an idea algebra up to
isomorphism.
\end{proof}

Murphy and Medin's empirical observations --- transfer from theories,
projection of unobserved features, the asymmetry between basic-level and
superordinate categories --- are then asymmetries of $\rho_T$ on $\Ideas_T$
together with the absence of any non-trivial $\eta$ on cross-domain inputs.

\formalstatus Theory-theory is the doctrine that conceptual processing is
local to a sub-algebra; \cref{cor:cogsci-theory} (a corollary of Theorems~1
and~8) makes the locality precise.

\litcomment The formalization does interesting work for theory-theory.
First, it forces an explicit representation of what was usually left
intuitive: a domain theory must designate the sub-monoid $\Ideas_T$ that
is closed under composition, the elevation operator $\eta_T$ that
governs which compositions produce theory-internal novelty, and the axiom
set $\Phi_T$ that anchors coherent compositions. Second, it supplies a
sharp operational definition of conceptual coherence: the composition
$a\compose b$ is coherent for $T$ iff its elevation lies in
$\langle\Phi_T\rangle_{\compose}$. Third, it makes cross-domain mismatch
the predicted generic case: the elevation $\eta(a,b)$ for $a\in\Ideas_T$
and $b\in\Ideas_{T'}$ generically exits both $\langle\Phi_T\rangle$ and
$\langle\Phi_{T'}\rangle$, and the felt incoherence of, e.g., a colorless
green idea is the cocycle obstruction $\omega(a,b)$ taking a value the
listener has no theory-internal interpretation for. The cognitive
architecture is therefore not a single monolithic algebra but a family of
sub-algebras with overlap, and the puzzle of how children acquire
theories \citep{MurphyMedin1985} becomes the algebraic puzzle of how the
sub-monoid $\Ideas_T$ is bootstrapped from a few generators and a small
set of axioms.

% =====================================================================
\section{Barsalou: perceptual symbols and ad-hoc categories}\label{sec:cogsci-barsalou}

\textbf{Original claim.} \citet{Barsalou1999} argued that concepts are
\emph{perceptual symbol systems}: simulators built from sensorimotor
records, deployed dynamically. Earlier, \citet{Barsalou1983} showed that
people fluently form ``ad-hoc categories'' such as \textsc{things to take
from a burning house} that have no stable mental representation but are
constructed on the fly to serve a goal.

\formalclaim Ad-hoc categories are \emph{transient quotient classes}.
Given a goal-context idea $g\in\Ideas$ (e.g., \textsc{burning house
escape}), define an equivalence on $\Ideas$ by
\[
  a\sim_g b \quad\iff\quad
  \rho(\conj{g}{a},\unit)=\rho(\conj{g}{b},\unit),
\]
the equivalence ``equally salient when conjugated by the goal context.''
The ad-hoc category $[a]_g$ is the equivalence class. Perceptual symbols
are the generators of a sub-monoid $\Ideas_P$ in which $\rho$ is realized
by sensorimotor similarity rather than by lexical co-occurrence; the
simulator that runs a perceptual symbol $p$ on context $c$ is
$\eta(p,c)$.

\begin{corollary}[Conjugation-coherence of ad-hoc
categories]\label{cor:cogsci-barsalou}
For any invertible $g\in\Ideas$, the relation $\sim_g$ is an equivalence
on $\Ideas$ and its quotient inherits the bilinear pairing $\rho_g$ given
by $\rho_g([a],[b])=\rho(\conj{g}{a},\conj{g}{b})$, well-defined modulo
$\ker\rho$.
\end{corollary}

\begin{proof}
Reflexivity, symmetry, and transitivity follow from those of equality on
$\RR$. Theorem~6 (Conjugation) makes $a\mapsto\conj{g}{a}$ an inner
automorphism preserving $\rho$, so the well-definedness of $\rho_g$ on
classes follows by the standard quotient argument modulo $\ker\rho$,
which is itself preserved by conjugation (Theorem~6).
\end{proof}

Thus ad-hoc categories are not a counterexample to a stable concept theory
but a feature of the fact that conjugation moves the resonance landscape
without altering the underlying algebra; the instability is the
context-dependence of $g$, not of the formalism.

\formalstatus Barsalou's perceptual symbols are sub-monoid generators;
ad-hoc categories are quotients by a context-conjugated resonance kernel,
guaranteed by Theorem~6 and \cref{cor:cogsci-barsalou}.

\litcomment The formalism gives Barsalou's two empirical pillars a common
algebraic root. Perceptual symbols are not a separate kind of mental
entity; they are simply the generators of a particular sub-monoid
$\Ideas_P$ on which $\rho$ is realized as sensorimotor similarity. A
``simulator'' for a perceptual symbol $p$ acting in context $c$ is just
the elevation $\eta(p,c)$, and the empirical observation that simulators
selectively activate modality-specific cortex is the prediction that
$\Ideas_P$ embeds non-trivially into the larger algebra in a way that
preserves the modality-tagged grading $\deg$. Ad-hoc categories, in
turn, are not strange exceptions: they are the predicted equivalence
classes whenever an idea $g$ acts as an inner automorphism on $\Ideas$
and one quotients by the kernel of the conjugated resonance. Both
phenomena --- the modality-specific simulator and the on-the-fly
category --- are consequences of Theorem~6, not additions to it. The
empirical question for the cognitive scientist is not whether such
structures exist (the algebra guarantees they do), but which generators
$p\in\Ideas_P$ a particular agent has acquired and which goal-context
ideas $g$ are habitually available.

% =====================================================================
\section{Fodor: Language of Thought}\label{sec:cogsci-lot}

\textbf{Original claim.} \citet{Fodor1975} argued that thinking requires a
combinatorial, syntactically structured medium of representation: a
\emph{Language of Thought} (LOT). Mental atoms are concept-words,
combination is governed by formal rules, and the productivity and
systematicity of thought are explained by the freeness of the resulting
structure.

\formalclaim Identify the LOT atoms with a generating set
$G\subseteq\Ideas$ such that the inclusion $G\hookrightarrow\Ideas$ extends
to an isomorphism
\[
  \Phi:\;\mathsf{FreeMon}(G)\;\xrightarrow{\;\cong\;}\;\langle G\rangle_{\compose},
\]
where $\mathsf{FreeMon}(G)$ is the free monoid on $G$. \emph{Productivity}
is the assertion that $\Phi$ is surjective onto a thinker's actually
entertainable ideas; \emph{compositionality} is the assertion that $\Phi$
is well-defined (the meaning of a concatenation is a function of its
parts); \emph{systematicity} is the further assertion (treated in
\cref{sec:cogsci-systematicity}) that $\rho$ respects the bilinear
structure inherited from the free vector space $V(G)$.

\begin{corollary}[Atomicity is a freeness
condition]\label{cor:cogsci-lot}
A sub-monoid $\Ideas_L\subseteq\Ideas$ admits a LOT presentation in the
sense of $\Phi$ iff every element of $\Ideas_L$ has a unique factorization
into $G$-atoms up to $\unit$. In that case $\omega|_{\Ideas_L}\equiv 0$
on the cohomological abelian group $A=\ZZ\langle G\rangle$.
\end{corollary}

\begin{proof}
Existence and uniqueness of factorization is the universal property of
$\mathsf{FreeMon}(G)$. Theorem~4 then gives $\omega(a,b)=v(a\compose b)
- v(a)-v(b)$ where $v$ counts atoms; this is identically zero on a free
monoid since the count of atoms is additive over concatenation.
\end{proof}

Equivalently, LOT is the limit case in which all conceptual emergence is
syntactic and none is semantic: $\eta=0$ at the level of atoms.

\formalstatus Fodor's atomism is the freeness of the generated sub-monoid;
\cref{cor:cogsci-lot}, a corollary of Theorems~1 and~4, gives the precise
algebraic content.

\litcomment The formalism vindicates the LOT hypothesis on its own terms
while exposing the cost of its purity. Vindication: a free sub-monoid
on a set of atoms has unique factorization, vanishing cocycle on the
atomic count, and full systematicity (\cref{cor:cogsci-fp}). Cost: the
freeness assumption forces $\eta$ to be zero on atomic compositions,
which prohibits semantic emergence at the lexical level. Most of what
the post-Fodorian literature on concepts complained about --- polysemy,
lexical creativity, dead and live metaphor, the productive coining of
ad-hoc compounds like \emph{teabag democracy} --- is the empirical
demand that $\eta$ be non-zero. The reconciliation supplied by the
algebra is straightforward: one is free to choose which sub-monoid is
LOT-like (free, with $\eta=0$) and which is not. Logical concepts and
explicit reasoning likely live in the former; lexical semantics and
discourse pragmatics in the latter. The disagreement between Fodor and
his critics, on this reading, is a disagreement about the location of
the boundary between the two regimes, not about whether either exists.

% =====================================================================
\section{Fodor \& Pylyshyn: systematicity}\label{sec:cogsci-systematicity}

\textbf{Original claim.} \citet{FodorPylyshyn1988} argued that any agent
that can think \textsc{John loves Mary} can also think \textsc{Mary loves
John}; a connectionist system that fails this generalization fails to
explain cognition. Systematicity, productivity, and inferential coherence
form an inseparable triad that cries out for classical
constituent-structure.

\formalclaim Let $G$ be a LOT generating set as in
\cref{sec:cogsci-lot}, with a binding sub-monoid $B\subseteq G$ of
\emph{relational} atoms. Bilinearity of $\rho$ on $V(\Ideas)$ delivers
systematicity automatically: for $r\in B$ and arguments $a_1,a_2$,
\[
  \rho\bigl(r\compose(a_1\compose a_2),\;r\compose(a_2\compose a_1)\bigr)
  \;=\;\rho(r,r)\cdot\rho(a_1\compose a_2,\,a_2\compose a_1),
\]
provided $\rho$ factors through $\compose$ on the relational layer (the
Smolensky tensor realization, \cref{sec:cogsci-smolensky}, gives a concrete
construction). The capacity to entertain one bound expression entails the
capacity to entertain its argument-permutation, since the bilinear form
treats both as definite vectors in the same finite-dimensional subspace.

\begin{corollary}[Systematicity from bilinearity]\label{cor:cogsci-fp}
If $\rho$ extends bilinearly (Theorem~2) and the relational layer is a
free sub-monoid (Theorem~1 plus \cref{cor:cogsci-lot}), then for every
$r,s\in B$ and every $a_1,\dots,a_k$ in the argument layer, every
permutation $r(a_{\pi(1)},\dots,a_{\pi(k)})$ is an element of $\Ideas$
with the same dimension and computable resonance.
\end{corollary}

\begin{proof}
The free sub-monoid contains all monomials in the generators; bilinearity
then makes the resonance with any other monomial well-defined. The
argument-permutation is itself a monomial. \qedhere
\end{proof}

Thus the systematicity argument, far from refuting connectionist or
algebraic semantics, becomes a \emph{theorem} once the formalism includes a
bilinear $\rho$; the empirical question is which monomials are actually
entertained, not which are in principle representable.

\formalstatus Systematicity is \cref{cor:cogsci-fp}, a direct corollary of
Theorems~1 and~2.

\litcomment The systematicity argument was meant as a knock-down case
against connectionism: if mental representations are merely vector
activations, what guarantees that thinking \textsc{John loves Mary}
suffices for thinking \textsc{Mary loves John}? The algebraic answer is
that bilinearity does. A connectionist system whose representations
inhabit a vector space and whose binding operation is bilinear (tensor
product, circular convolution, multi-head attention) automatically
satisfies systematicity for the represented relations. Modern
distributed systems, including transformers, satisfy bilinearity by
construction; the empirical question is therefore not whether
systematicity \emph{can} hold in such systems but whether the training
distribution exercises the relevant permutations sufficiently for the
network to learn them. \citet{FodorPylyshyn1988} were therefore right
about the constraint and wrong about the architecture: the constraint
is a theorem of the algebra, and any architecture that realizes the
algebra automatically satisfies it. The post-1988 literature has, to a
considerable extent, vindicated this reading; the persistence of the
debate reflects ongoing disagreement about whether trained networks
actually realize the bilinear structure or only approximate it.

% =====================================================================
\section{Fauconnier \& Turner: conceptual blending}\label{sec:cogsci-blending}

\textbf{Original claim.} \citet{FauconnierTurner2002} hold that human
imagination routinely constructs \emph{blended mental spaces} in which
selected structure from two or more input spaces is projected into a novel
space that is neither input but acquires \emph{emergent structure}. A
``land yacht'' is a car that inherits land travel and yacht-like luxury,
producing emergent inferences (no sails, no water) absent from the inputs.

\formalclaim A mental space is an idea $a\in\Ideas$. A blend of inputs
$a,b\in\Ideas$ is exactly the elevation component of their composition,
\[
  \mathrm{Blend}(a,b)\;\eqdef\;\eta(a,b),
\]
with the preserved structure $\sigma(a,b)$ playing the role of the
``generic space'' and the negated structure $\nu(a,b)$ playing the role of
clashes that the blend resolves. The Aufhebung Decomposition Theorem
(Theorem~3) is, read in cognitive-scientific terms, the existence and
uniqueness of the blend up to resonance equivalence.

\begin{corollary}[Blend uniqueness]\label{cor:cogsci-blend}
For any inputs $a,b\in\Ideas$ the blend $\eta(a,b)$ is unique modulo
$\ker\rho$, so two blends of the same inputs differ only by a
resonance-invisible kernel element. Iterated blends form a 2-cocycle
chain whose obstruction is exactly $\omega(a,b)$.
\end{corollary}

\begin{proof}
Uniqueness modulo $\ker\rho$ is Theorem~3. Iterating
$\mathrm{Blend}(\mathrm{Blend}(a,b),c)$ versus
$\mathrm{Blend}(a,\mathrm{Blend}(b,c))$ yields a discrepancy controlled by
$\omega(a,b\compose c)-\omega(a\compose b,c)+\omega(b,c)-\omega(a,b)$,
which vanishes by the cocycle identity (Theorem~4). \qedhere
\end{proof}

The graded family of mental spaces is then the filtration of $\Ideas$ by
$\deg$ (Theorem~9): low-degree spaces are local, concrete, and quickly
composed; high-degree spaces accumulate emergent content via
$\omega$-additions.

\formalstatus Blending is the elevation operator $\eta$ supplied by
Theorem~3; emergent inferences are the cocycle $\omega$ of Theorem~4;
\cref{cor:cogsci-blend} packages these into the existence and
chain-coherence of blends.

\litcomment The Fauconnier--Turner programme has been criticized for
proliferation: every novel inference becomes a ``blend'' and the theory
risks unfalsifiability. The algebra disciplines the proliferation. Not
every composition is a blend; only those whose elevation $\eta(a,b)$
contributes content not already in $\sigma(a,b)$ counts. The cocycle
$\omega(a,b)$ measures exactly how much new content is added, and the
chain identity of \cref{cor:cogsci-blend} guarantees that iterated blends
do not accumulate spurious structure: any apparent gain from re-grouping
is recoverable from the cocycle. A predicted experimental signature
follows: if blending is genuinely $\eta$-driven, then the cognitive
processing time and ERP amplitude associated with comprehending a
novel blend should track the magnitude of $\omega(a,b)$ in some
behavioural norming, not the surface lexical novelty alone. To the
extent that this prediction has been confirmed (in the experimental
work on counterfactuals, deontic blends, and arithmetic blends), the
algebra ratifies the original programme; to the extent it has not, the
algebra circumscribes the programme to the cases where $\omega$ is
non-trivial.

% =====================================================================
\section{Lakoff: conceptual metaphor}\label{sec:cogsci-lakoff}

\textbf{Original claim.} \citet{LakoffJohnson1980} argued that abstract
concepts are systematically structured by metaphorical mappings from
concrete source domains: \textsc{argument is war}, \textsc{time is
money}, \textsc{theories are buildings}. The mappings preserve image
schemas (containers, paths, force dynamics) and yield productive new
inferences in the target.

\formalclaim Let $S,T\subseteq\Ideas$ be source and target sub-algebras
(\cref{cor:cogsci-theory}). A \emph{conceptual metaphor} is a partial
graded monoid morphism $M:S\rightharpoonup T$ such that
\begin{equation}\label{eq:cogsci-lakoff-rho}
  \rho_T(M(a),M(b))\;=\;\lambda_M\,\rho_S(a,b)
  \qquad\text{for some }\lambda_M\in\RR_{>0},
\end{equation}
on the domain of $M$. The metaphor is \emph{productive} when $M$ commutes
with $\compose$ in the sense $M(a\compose b)=M(a)\compose M(b)$ wherever
defined. Image schemas are the $\sigma$-invariants of $M$:
$\sigma_T(M(a),M(b)) = M(\sigma_S(a,b))$.

\begin{corollary}[Cross-domain $\rho$ as systematic
mapping]\label{cor:cogsci-lakoff}
Equation~\eqref{eq:cogsci-lakoff-rho} together with the morphism property
of $M$ implies that $M$ extends uniquely to a graded vector-space
homomorphism $V(S)\to V(T)$ that is a similarity, hence preserves the
bilinear shadow of all higher-order structure (the cocycle $\omega$
transports as $\omega_T(M(a),M(b))=M_*\omega_S(a,b)$ on the abelian
extension).
\end{corollary}

\begin{proof}
Bilinear extension (Theorem~2) and uniqueness of the Aufhebung
decomposition modulo $\ker\rho$ (Theorem~3) make any morphism preserving
$\rho$ up to scale automatically preserve $\sigma$ on inputs in its
domain. Theorem~4's normalized 2-cocycle property is preserved by any
monoid map that intertwines the valuations.
\end{proof}

Failures and contestations of conceptual metaphor (Murphy's critique;
asymmetry of source and target) are the failure of $M$ to extend beyond a
small generating set: $\lambda_M$ depends on the sub-domain, and $\rho$'s
intrinsic non-symmetry (curvature $\kappa$) prevents the metaphor from
running in both directions. The standard observation that one can say
\emph{her argument collapsed} but not \emph{her building was refuted} is,
on this reading, the report that $M:\mathit{building}\to\mathit{argument}$
preserves $\rho$ in one direction only; the inverse mapping fails the
$\rho$-preservation condition because the relevant antisymmetric part of
$\rho$ (the curvature) is non-zero, and curvature blocks the construction
of a two-sided morphism. The asymmetry of metaphor is therefore an
algebraic fact, not a contingent linguistic convention.

\formalstatus Conceptual metaphor is a partial graded morphism between
sub-algebras preserving $\rho$ up to scale; \cref{cor:cogsci-lakoff} is
its corollary status from Theorems 2, 3, and 9.

\litcomment The Lakoff--Johnson programme has been criticized for the
profligate generation of metaphor schemata; every linguistic figure
becomes evidence for some cross-domain mapping. The algebra supplies a
test: a candidate mapping $M$ qualifies as a conceptual metaphor only
when it is partial morphism preserving $\rho$ up to scale. This
discipline excludes mere lexical analogy and admits only mappings whose
preservation of $\rho$ produces systematic novel inferences in the
target domain. The image-schema doctrine receives an even cleaner
algebraic reading: image schemas are exactly the preserved
$\sigma$-structure of $M$, that is, the parts of the source that the
metaphor allows to survive into the target. Containment, path, and
balance schemas are then $\sigma$-invariants of broad classes of
metaphors, which explains their cross-cultural recurrence: any partial
morphism preserving $\rho$ on the spatial sub-algebra will preserve
these structural invariants, regardless of the source or target.

% =====================================================================
\section{Fillmore: frame semantics}\label{sec:cogsci-fillmore}

\textbf{Original claim.} \citet{Fillmore1982} argued that lexical meaning
is anchored in \emph{frames}: structured background scenes that supply
roles, fillers, and inferences. To understand \textsc{buy} is to access
the commercial transaction frame, which involves \textsc{buyer},
\textsc{seller}, \textsc{goods}, \textsc{money}. Lexical units evoke the
frame; \emph{profiling} selects which role is foregrounded.

\formalclaim A \emph{frame} is a sub-monoid $\Ideas_F=\langle
L_F\rangle_{\compose}$ generated by a finite set $L_F\subseteq\Ideas$ of
lexical units, equipped with a finite role set
$R_F=\{r_1,\dots,r_n\}\subseteq\Ideas$ and a set of role-binding axioms
$r_i\compose f = f\compose r_i$ for $f\in L_F$ (i.e., the roles commute
with the lexical units they apply to). \emph{Profiling} is the
$\sigma$-projection: profiling lexical unit $\ell$ in frame $F$ produces
$\sigma(\ell,\Phi_F)$ where $\Phi_F=\bigodot_{i}r_i$ is the iterated
composition of the frame's roles. Different lexical entries can profile
the same generic scene with different $\sigma$-residue (the buyer/seller
asymmetry).

\begin{corollary}[Profiling as Aufhebung
projection]\label{cor:cogsci-frames}
For each $\ell\in L_F$ the profiling $\sigma(\ell,\Phi_F)$ is uniquely
defined modulo $\ker\rho$ and varies with $\ell$ within a single frame.
The grading $\deg$ restricts to $\Ideas_F$ and assigns each frame element
a degree summing the degrees of its constituent roles and lexical units.
\end{corollary}

\begin{proof}
Theorem~3 supplies uniqueness of $\sigma$ modulo $\ker\rho$. Theorem~9
extends $\deg$ as a $G$-grading compatible with $\compose$ on every
sub-monoid; commuting roles ensure additivity on the frame. \qedhere
\end{proof}

Frame inheritance \citep{Fillmore1982} corresponds to nesting of
sub-monoids $\Ideas_{F'}\subseteq\Ideas_F$ with their associated roles,
and frame mismatches register as nonzero $\nu$-components when one frame's
output is fed into another's input slot.

\formalstatus A frame is a sub-monoid (Theorem~1); profiling is
$\sigma$-projection (Theorem~3); the degree of a profiled element is
extracted via Theorem~9 --- packaged as \cref{cor:cogsci-frames}.

\litcomment Frame semantics, in its FrameNet implementation, has long
suffered from the appearance of being a vast hand-curated database
rather than a theory. The algebra makes the theoretical content explicit:
each frame is a sub-monoid with a designated set of role generators, and
the difference between two near-synonyms (e.g.\ \emph{buy} and
\emph{sell}) is captured by the differing $\sigma$-projection onto the
shared role structure. Frame inheritance becomes literal sub-monoid
inclusion, and the partial-order relations FrameNet records (Inheritance,
Subframe, Perspective\_on, Causative\_of) become specific structural
relations between sub-monoids. This recasting also explains a common
puzzle: why frame elements that are obligatory in one frame become
optional or inferred in an inherited frame. The answer is that the
inherited sub-monoid retains the role generator but does not require its
non-trivial profiling; $\sigma(\ell,r)$ may be zero for some inheriting
$\ell$, so the role is present in principle but invisible to behavior.

% =====================================================================
\section{G\"ardenfors: conceptual spaces}\label{sec:cogsci-gardenfors}

\textbf{Original claim.} \citet{Gardenfors2000} proposed that concepts are
\emph{convex regions} in a multidimensional \emph{conceptual space} whose
quality dimensions are anchored in perception (color in HSL, pitch on a
log-frequency line, taste on the tetrahedron of basic tastes). Categories
are convex; prototypes are centroids; similarity is inverse distance.

\formalclaim Pass to the free $\RR$-vector space $V(\Ideas)$. A
\emph{quality dimension} is a $\rho$-orthogonal one-dimensional subspace
$D_i\subseteq V(\Ideas)$. A \emph{conceptual space} of dimension $n$ is
the orthogonal direct sum $V_n=\bigoplus_{i=1}^{n}D_i$ with the
restriction of $\rho$ to $V_n$ assumed positive-definite (this constrains
the choice of dimensions). A \emph{concept} is a convex subset
$C\subseteq V_n$. The \emph{$\rho$-induced metric} is
\[
  d_\rho(x,y)\;\eqdef\;\sqrt{\rho(x-y,\,x-y)},
\]
which is a genuine metric by positive-definiteness of $\rho|_{V_n}$.
Betweenness is then the standard metric betweenness, and prototypes are
the $d_\rho$-centroids of $C$.

\begin{corollary}[Convex regions are stable under
composition-respecting reweighting]\label{cor:cogsci-gardenfors}
Let $C\subseteq V_n$ be convex, $g\in\Ideas$ invertible. Then
$\conj{g}{C}\eqdef\{\conj{g}{x}:x\in C\}$ is convex in
$\conj{g}{V_n}$, and the prototype satisfies
$\pi_{\conj{g}{C}}=\conj{g}{\pi_C}$.
\end{corollary}

\begin{proof}
Conjugation by $g$ is a linear isomorphism on $V(\Ideas)$ (Theorem~6) and
linear maps preserve convexity and centroids. \qedhere
\end{proof}

Voronoi tessellations of $V_n$ around a finite set of prototype points
generate convex categories \citep{Gardenfors2000}, and these tessellations
are equivariant under the conjugation action; the algebra is therefore
compatible with G\"ardenfors's geometry while extending it to a
non-commutative $\rho$-curvature regime that conceptual-space theory does
not natively address.

\formalstatus G\"ardenfors's convex-region thesis is the
positive-definite, orthogonal-direct-sum specialization of $\rho$ on
$V(\Ideas)$, with conjugation invariance from \cref{cor:cogsci-gardenfors}
(Theorems~2 and 6).

\litcomment G\"ardenfors's geometric vision of cognition has the virtue
of making cognitive theory metric and the limitation of restricting
itself to symmetric, positive-definite similarity. The algebra extends
the geometry to the non-symmetric and non-positive cases that
G\"ardenfors's setup excludes. Where conceptual-space theory must locate
asymmetric similarity outside the geometry (and treat it as an
exception), the algebra puts asymmetry inside, in the curvature
$\kappa$. Where conceptual-space theory must restrict itself to
quality dimensions that admit a meaningful inner product, the algebra
admits arbitrary $\rho$ and recovers the inner-product case as the
positive-definite restriction. The cognitive science of color, taste,
and pitch (where conceptual-space theory excels) inhabits the
restricted regime; the cognitive science of moral, social, and
narrative concepts inhabits the curved regime, where asymmetric
similarity is the rule rather than the exception.

% =====================================================================
\section{Smolensky: tensor product variable binding}\label{sec:cogsci-smolensky}

\textbf{Original claim.} \citet{Smolensky1990} introduced \emph{tensor
product representations} as a means of giving connectionist networks a
combinatorial syntax: a structured object is a sum of role/filler tensor
products, $\sum_i f_i\otimes r_i$, with binding by $\otimes$ and unbinding
by inner product with the role vector.

\formalclaim Realize $\Ideas$ as a sub-monoid of the tensor algebra
$T(V)=\bigoplus_{k\ge 0}V^{\otimes k}$ over a vector space $V$ of fillers
and roles. Set
\[
  a\compose b \;\eqdef\; a\otimes b,
\]
with $\unit$ the scalar $1\in V^{\otimes 0}$. Resonance is the natural
inner product:
\[
  \rho(a,b)\;\eqdef\;\langle a,b\rangle_{T(V)}
  \;=\;\sum_k \langle a_k,b_k\rangle_{V^{\otimes k}}.
\]
The Aufhebung decomposition reads off as
$\sigma(a,b)=$ the symmetrization of $a\otimes b$, $\nu(a,b)=$ the
antisymmetrization, $\eta(a,b)=$ the irreducible higher-tensor remainder.
$\omega(a,b)$ records the rank increment in the grading
$\deg(a)=$ tensor rank.

\begin{corollary}[Tensor algebra is a model of the idea
algebra]\label{cor:cogsci-smol}
The tuple $(T(V),\otimes,1,\langle\cdot,\cdot\rangle,\sigma,\nu,\eta,
\omega,\mathrm{rank})$ satisfies the foundational axioms; in particular
$\omega$ is the rank-difference cocycle and is normalized.
\end{corollary}

\begin{proof}
Associativity of $\otimes$ and the unit law for $1$ are standard.
Bilinearity of $\langle\cdot,\cdot\rangle$ on each $V^{\otimes k}$
extends to $T(V)$ by orthogonal direct sum. The Aufhebung pieces are
the standard sym/antisym/remainder decomposition. The cocycle
$\omega(a,b)=\mathrm{rank}(a\otimes b)-\mathrm{rank}(a)-\mathrm{rank}(b)
=0$ on pure tensors, recovering normalization. \qedhere
\end{proof}

Smolensky's binding is therefore a \emph{model} of the idea algebra: the
formalism does not contradict it but realizes it as one option among
others. The cost (exponential growth of dimension with depth) is exactly
the predicted cost of refusing $\omega$ to take values outside the trivial
abelian group.

\formalstatus The tensor algebra $T(V)$ is an explicit model of
$\IdeaAlg$; existence is \cref{cor:cogsci-smol}, a transparent application
of Theorems~1--4.

\litcomment The Smolensky construction occupies a privileged position in
the algebra: it is the unique up-to-isomorphism faithful realization in
which all of compositional structure is encoded in tensor rank, every
binding is invertible, and every cocycle vanishes. This faithfulness is
purchased at the cost of unbounded dimensionality; cognition cannot
literally implement $T(V)$ because the brain is finite. The Plate HRR
construction (\cref{sec:cogsci-plate}) is then the natural bandwidth-
limited compromise. The two together delineate a tradeoff curve: any
realization of the idea algebra must choose where on the rank--noise
plane it sits, and the cognitive system seems to make this choice
adaptively, increasing rank for conscious deliberate composition and
compressing into noisy fixed-dimension form for fast, automatic
processing. The algebra makes this tradeoff a theorem rather than an
engineering choice.

% =====================================================================
\section{Plate: holographic reduced representations}\label{sec:cogsci-plate}

\textbf{Original claim.} \citet{Plate1995} developed Holographic Reduced
Representations (HRRs) to address the dimensional explosion of tensor
products: bind by \emph{circular convolution} $\circledast$ and unbind by
circular correlation. Dimension is preserved; binding is approximate but
rich; cleanup memory recovers exact structure.

\formalclaim Take $V=\RR^n$ for fixed $n$, and define $\compose$ on
$\Ideas\eqdef\RR^n$ by
\[
  (a\compose b)_k\;\eqdef\;(a\circledast b)_k
  \;=\;\sum_{j=0}^{n-1}a_j\,b_{(k-j)\bmod n}.
\]
Identify $\unit$ with the impulse $e_0=(1,0,\dots,0)$. Resonance is the
inner product $\rho(a,b)=\langle a,b\rangle$. Crucially, for random
unit-norm vectors $a,b,c,d$,
\[
  \mathbb{E}\!\bigl[\rho(a\circledast b,\,c\circledast d)\bigr]
  \;=\;\rho(a,c)\,\rho(b,d) + O(1/n),
\]
i.e., $\rho$ is approximately preserved under composition. This is the
\emph{compositional resonance} condition.

\begin{corollary}[$\rho$-preserving compositional
operator]\label{cor:cogsci-plate}
Circular convolution gives a (commutative, associative) realization of
$\compose$ on $\RR^n$ for which $\rho$ is asymptotically preserved with
$O(1/\sqrt{n})$ noise. The cocycle $\omega$ is the residual
$\rho(a\compose b,\,c\compose d)-\rho(a,c)\rho(b,d)$, valued in $\RR$.
\end{corollary}

\begin{proof}
Commutativity and associativity of $\circledast$ are standard
properties of circular convolution. The expected-value identity is the
classical HRR result \citep{Plate1995} from independence of random
phases. The cocycle interpretation follows from Theorem~4 by reading
the residual as the obstruction to multiplicativity of $\rho$. \qedhere
\end{proof}

HRRs thus realize a non-faithful but bandwidth-bounded model of the idea
algebra: the cost of compression appears as cocycle noise, predicted by
the formalism.

\formalstatus Plate's HRR is a fixed-dimension realization of $\compose$
and $\rho$ where compositional preservation is approximate;
\cref{cor:cogsci-plate} (corollary of Theorems~2 and~4) names the cocycle
that records the approximation error.

\litcomment HRR illustrates a general principle: any bandwidth-bounded
realization of the idea algebra must accept a non-trivial cocycle. The
$1/\sqrt{n}$ noise of circular convolution is not a defect of HRR
specifically but the price of compression: information-theoretically,
binding $k$ atoms into a $d$-dimensional vector with the precision
required for unique unbinding requires $d=\Omega(k)$ for exact
preservation, and any sublinear $d$ pays in cocycle noise. The
algebraic prediction is that biological cognition --- itself
bandwidth-bounded --- should exhibit a structurally similar noise
floor on bound-representation accuracy, and that improvements in
working memory should appear as increases in effective $d$ rather than
as algorithmic improvements. The empirical literature on working
memory capacity is broadly consistent with this prediction.

% =====================================================================
\section{Mikolov / GloVe: word embeddings}\label{sec:cogsci-mikolov}

\textbf{Original claim.} \citet{Mikolov2013} and
\citet{PenningtonGloVe2014} demonstrated that word vectors learned from
large corpora encode \emph{semantic similarity} in cosine angles and
\emph{semantic relations} in linear differences:
\[
  \mathbf{v}(\mathit{king})-\mathbf{v}(\mathit{man})+
  \mathbf{v}(\mathit{woman})\;\approx\;\mathbf{v}(\mathit{queen}).
\]
This was the first persuasive evidence that distributional learning yields
algebraic structure.

\formalclaim Let $V_W\subseteq V(\Ideas)$ be the linear subspace spanned
by a vocabulary $W\subseteq\Ideas$ embedded by a learned map
$\mathbf{v}:W\to\RR^d$. Treat the symmetric part of $\rho$ on $V_W$ as
the cosine-induced inner product
$\rho_{\mathrm{sym}}(a,b)=\langle\mathbf{v}(a),\mathbf{v}(b)\rangle/
(\|\mathbf{v}(a)\|\,\|\mathbf{v}(b)\|)$. The analogy
$\mathit{king}-\mathit{man}+\mathit{woman}=\mathit{queen}$ is then the
statement that the gender relation acts by a \emph{conjugation} in the
free vector space:
\[
  \mathbf{v}(\mathit{queen})\;=\;\mathbf{v}(\mathit{king})+
  \bigl(\mathbf{v}(\mathit{woman})-\mathbf{v}(\mathit{man})\bigr),
\]
which is the linearization of the inner-automorphism action
$\conj{g}{(\,\cdot\,)}$ for $g$ chosen so that
$\mathbf{v}(g)-\mathbf{v}(g^{-1})=\mathbf{v}(\mathit{woman})-\mathbf{v}(\mathit{man})$.

\begin{corollary}[Analogy as conjugation in
$V(\Ideas)$]\label{cor:cogsci-mikolov}
On the subspace where $\rho$ is symmetric and positive-definite,
conjugation by an invertible idea $g$ acts as a translation in the
linearization $V_W$. Word-embedding analogy ``$a:b::c:d$'' is then the
identity $\mathbf{v}(d)=\mathbf{v}(c)+\mathbf{v}(b)-\mathbf{v}(a)$,
recoverable as the projection of an exact conjugation onto $V_W$.
\end{corollary}

\begin{proof}
By Theorem~6, conjugation by an invertible $g$ is an inner automorphism
preserving $\rho$. Linearizing the multiplicative group action on
$V(\Ideas)$ around the identity yields an additive Lie-algebra action,
i.e.\ translation. Restricting to $V_W$ gives the embedding identity.
\qedhere
\end{proof}

Failures of the analogy (regularization-induced errors, polysemy
collapse, gender bias) reflect the fact that $\rho$ is not symmetric in
general (curvature $\kappa\ne 0$): the linearization is only a first-order
approximation to the true non-commutative structure.

\formalstatus Cosine similarity is a symmetric specialization of $\rho$;
analogy is the linear shadow of conjugation (Theorem~6); the analogy
identity is \cref{cor:cogsci-mikolov}.

\litcomment The discovery of analogical structure in word embeddings
\citep{MikolovKingQueen2013} was striking precisely because
distributional learning was widely thought to be too crude to encode
algebraic relations. The algebra removes the surprise. A learning
process that minimizes a quadratic loss on a co-occurrence matrix is
implicitly fitting a symmetric bilinear form on $V_W$; the bilinear
form is a $\rho$-restriction; and conjugation by an invertible idea
linearizes to a translation in any $\rho$-respecting embedding. The
``king $-$ man $+$ woman $=$ queen'' identity is therefore not a
miracle but the predicted shadow of Theorem~6 in any vocabulary whose
relational structure has been faithfully captured. The well-known
failures of the analogy --- bias along stereotyped axes, failure on
infrequent words, breakdown under polysemy --- are the predicted
breakdowns when the learned $\rho$ is non-symmetric and the
linearization assumption fails. The algebra thus simultaneously
explains the success and the failure modes of distributional analogy.

% =====================================================================
\section{Transformers as latent ideas}\label{sec:cogsci-transformers}

\textbf{Original claim.} \citet{Vaswani2017} introduced the Transformer
architecture, in which token representations are updated by repeated
\emph{attention} blocks: each token $x_i$ produces query, key, and value
vectors $q_i,k_i,v_i$, and the new representation is
$x_i'=\sum_j \mathrm{softmax}_j(q_i^\top k_j/\sqrt{d})\,v_j$. In-context
learning \citep{BrownGPT3} is the empirical observation that, given a few
examples in the prompt, the network executes a new task without parameter
updates.

\formalclaim The (un-normalized) attention map is a learned, asymmetric
bilinear form
\[
  \rho_{\mathrm{att}}(x_i,x_j)\;\eqdef\;\frac{q_i^\top k_j}{\sqrt{d}},
\]
i.e.\ a non-symmetric resonance pairing on the latent idea space spanned by
the token activations. The output update is the resonance-weighted
composition
\[
  x_i' \;=\; x_i\compose_{\mathrm{att}}\bigl(\textstyle\sum_j
  \mathrm{softmax}_j(\rho_{\mathrm{att}}(x_i,x_j))\,v_j\bigr).
\]
In-context learning is on-the-fly conjugation: the prompt prefix
$g_{\mathrm{ctx}}$ acts as an invertible context idea, and the attention
machinery computes $\conj{g_{\mathrm{ctx}}}{x}$ for each downstream token
$x$, transporting the resonance landscape to the task-appropriate frame
without altering the underlying parameters.

\begin{corollary}[Attention as learned non-symmetric
$\rho$]\label{cor:cogsci-transformers}
The attention bilinear form $\rho_{\mathrm{att}}$ is a non-symmetric
specialization of $\rho$; its antisymmetric part is a curvature
$\kappa_{\mathrm{att}}$ in the sense of Theorem~5, and the in-context
learning operation is the corresponding inner-automorphism action in the
sense of Theorem~6.
\end{corollary}

\begin{proof}
Bilinearity of $q_i^\top k_j$ in $x_i$ and $x_j$ (when $q$ and $k$ are
linear projections) verifies the bilinear form axioms of Theorem~2.
Asymmetry follows because $W_Q\ne W_K$ in general; the antisymmetric part
is a closed 2-cochain (Theorem~5). The conjugation interpretation of
in-context learning is Theorem~6 applied to the prompt prefix as
invertible idea (invertibility holding generically on the relevant
subspace). \qedhere
\end{proof}

The empirically observed phenomena --- induction heads, in-context
classification, curriculum effects --- are then specializations of the
formalism. The non-trivial cohomology class of $\kappa_{\mathrm{att}}$ is
the algebraic correlate of the irreducibility of attention to a symmetric
similarity.

\formalstatus Attention is a learned non-symmetric $\rho$; in-context
learning is conjugation; both are guaranteed by Theorems~2, 5, and 6,
packaged as \cref{cor:cogsci-transformers}.

\litcomment The transformer is, of all the architectures considered in
this chapter, the one whose alignment with the algebra is the closest.
The query/key projections produce a learned non-symmetric resonance
pairing; the value projection acts as the elevation operator; the
residual stream realizes composition; the multi-head structure
implements parallel restriction to multiple sub-monoids; layer
normalization rescales $\rho$ to keep it in a numerically tractable
regime; and in-context learning is a literal conjugation of the
representation by the prompt prefix. Each of these correspondences is
not metaphorical: it is structural, in the sense that the bilinear
forms involved satisfy the axioms of \cref{ch:foundations} verbatim.
The empirical success of transformers across linguistic and
non-linguistic domains is, on this reading, evidence that the idea
algebra is the natural target abstraction for any cognitive
architecture --- biological or artificial --- that must compose
high-dimensional structured representations under bandwidth and
latency constraints.

% =====================================================================
\section{A unifying picture}\label{sec:cogsci-unification}

The fifteen positions of this chapter span every layer of cognitive
architecture from sensorimotor records to large-scale neural language
models. They have, historically, been pitted against one another:
prototypes vs.\ definitions, simulators vs.\ symbols, classical vs.\
connectionist. The idea algebra reveals these oppositions to be artifacts
of partial views.

\begin{theorem}[Layered idea-algebra architecture]\label{thm:cogsci-layered}
Every position in
\cref{sec:cogsci-classical,sec:cogsci-rosch,sec:cogsci-theorytheory,
sec:cogsci-barsalou,sec:cogsci-lot,sec:cogsci-systematicity,
sec:cogsci-blending,sec:cogsci-lakoff,sec:cogsci-fillmore,
sec:cogsci-gardenfors,sec:cogsci-smolensky,sec:cogsci-plate,
sec:cogsci-mikolov,sec:cogsci-transformers} is the restriction of the
foundational idea algebra $\IdeaAlg$ to (i) a sub-monoid (a category, a
domain theory, a frame, a vocabulary, a layer of activations); (ii) a
choice of which operators are non-trivial ($\rho$ alone for
distributional models; $\sigma$ for image schemas; $\eta$ for blending);
and (iii) a regime of $\omega$ (zero for atomic LOT, controlled for HRR,
graded for transformers). Each position is the formalism viewed through
one of these three filters.
\end{theorem}

\begin{proof}[Sketch]
Each section above gave the explicit restriction. By Theorem~8
(Structural Equivalence), restricting to a sub-monoid and choosing the
non-trivial operators preserves idea-algebra structure up to
isomorphism; by Theorem~9, the grading is compatible with each
restriction. The third filter (regime of $\omega$) is the choice of
abelian group $A$ for cohomology, which Theorem~4 leaves free. The
union of the fifteen restrictions, glued along their shared sub-monoids,
returns the full algebra $\IdeaAlg$. \qedhere
\end{proof}

The picture is layered, not contested: \emph{perception} provides the
fillers (Barsalou); \emph{language} a generating set (Fodor, Fillmore);
\emph{categorization} a graded restriction of $\rho$ (Rosch, G\"ardenfors,
Mikolov); \emph{theory} a sub-algebra (Murphy \& Medin); \emph{metaphor}
a partial morphism (Lakoff); \emph{imagination} an $\eta$-operation
(Fauconnier \& Turner); \emph{binding} a $\compose$-realization
(Smolensky, Plate); \emph{contextual computation} a learned non-symmetric
$\rho$ (Vaswani, Brown). The cognitive architecture is the limit of the
diagram these layers form in the category of idea algebras. The next
chapter takes up the connectionist side of this story in its own right;
the present one closes by registering that the long-running war between
the symbolic and the statistical is, formally, a war between two valid
sub-restrictions of a single algebra.

\formalstatus The unifying claim is \cref{thm:cogsci-layered}, an
application of Theorems~1, 4, 8, and 9. The chapter as a whole
demonstrates that the idea algebra is not a competitor to existing theories
of concepts but a common substrate in which each is a coherent
specialization.

\subsection*{What the unification buys}

The translation effected in this chapter is not a mere notational
re-skinning. It buys at least four things that the cognitive science of
concepts has lacked since Rosch.

\emph{First, a precise vocabulary for partial agreement.} Theorists of
concepts typically agree on more than they disagree on, but lack the
vocabulary to state the agreement precisely. The classical and prototype
theories agree about the existence of features and disagree about how
membership in a category depends on those features. Translated, they
agree on the existence of the sub-monoid and disagree on the form of
$\rho$. Theory-theory and prototype theory agree about graded
membership and disagree about whether category coherence is statistical
or causal-explanatory. Translated, they agree about the bilinear
extension of $\rho$ and disagree about whether the relevant sub-algebra
is open under arbitrary composition or only under
theory-internal axioms. Such precise statements of agreement open
tractable empirical questions; impressionistic statements do not.

\emph{Second, a roadmap for hybrid models.} The empirical successes of
recent years (in word embeddings, in transformer-based language models,
in conceptual blending experiments) are evidence not for any one
position but for the algebraic substrate they all share. A hybrid
cognitive architecture that combines symbolic LOT-style atoms with
distributed transformer-style continuous resonance is, in the algebra,
just an architecture in which one sub-monoid is free (LOT) and another
admits a non-trivial cocycle (transformer). The two sub-monoids share
the same generators (lexical items) and the same composition operator
$\compose$, but differ in the regime of $\omega$. This is a
formally clean way to specify hybrid architectures without ad-hoc
glue code.

\emph{Third, falsifiable predictions.} Each section above includes a
formal status claim, and each formal status claim is falsifiable. If
$\rho$ is genuinely bilinear, then conjugation should preserve
typicality; if it does not, the bilinearity assumption fails locally
and the theory is wrong about the relevant sub-domain. If $\eta(a,b)$
is genuinely the cognitive substrate of conceptual blending, then the
processing time of a blend should track the magnitude of the
elevation, not the surface lexical novelty. If in-context learning is
genuinely conjugation, then the in-context behavior of a transformer
should be invariant under the inverse-conjugation of the prompt
prefix. Each of these is a non-trivial empirical prediction.

\emph{Fourth, mathematical economy.} The fifteen positions of this
chapter previously required fifteen distinct vocabularies, fifteen
distinct sets of constructs, fifteen distinct experimental traditions.
The algebra requires one vocabulary, nine theorems, and a finite list
of restrictions. The economy is real: anything provable in the unified
formalism is automatically true in any of the fifteen sub-models that
specialize it, and any empirical prediction made within one
sub-model can be tested against the data collected in any other
sub-model that shares the relevant sub-monoid.

\subsection*{Open questions}

Three open questions follow naturally from the unification.

The first concerns \emph{learning the algebra from data}. Modern
distributional models (word embeddings, transformers) demonstrate that
some of the algebra's structure can be acquired from co-occurrence
statistics. What is the minimal supervisory signal sufficient to
acquire the full algebra, including a non-trivial $\eta$ and
$\omega$? The conjecture, supported by the in-context learning
phenomena of \cref{sec:cogsci-transformers}, is that next-token
prediction over diverse text already constitutes a sufficient signal.
A precise answer would tighten the connection between the algebra
and the engineering of large language models.

The second concerns \emph{the neural substrate}. The algebra is, by
construction, hardware-neutral: it can be realized in symbols, in
tensors, in HRRs, in attention, or in spike timings. Which substrate
the brain uses is a separate empirical question. The algebra, however,
constrains the search: any biologically plausible substrate must
support a bilinear $\rho$, an elevation operator $\eta$, and a
non-trivial cocycle $\omega$. The phasic-burst codes of cortical
microcircuits, the temporal-coincidence codes of hippocampus, the
predictive coding of cortex --- each is a candidate, and each is
testable against the algebraic constraints.

The third concerns \emph{the limits of formalization}. The algebra
captures composition, resonance, Aufhebung, emergence, and grading.
It does not (yet) capture affect, motivation, embodied action, or
the social character of meaning. These are not absences in the
algebra so much as places where additional structure must be added.
Subsequent chapters take up several of these strands: \cref{ch:lakoff}
extends the algebra to image schemas; \cref{ch:vsa} elaborates the
distributed realization; \cref{ch:emergence} returns to the cocycle
and its cohomological invariants. The cognitive science of concepts,
in short, is not finished by this chapter; it is now equipped to
finish itself.

\subsection*{Coda}

The history of cognitive science can be read as a sequence of partial
discoveries about a single mathematical object. The classical theorists
discovered that some compositions are exactly Boolean. Rosch discovered
that resonance is graded. Murphy and Medin discovered that resonance is
local to sub-algebras. Fodor discovered that some sub-algebras are free.
Fauconnier and Turner discovered that elevation is a distinct operator
from preservation. Lakoff discovered that partial morphisms transport
resonance across sub-algebras. Smolensky and Plate discovered concrete
realizations of composition. Mikolov discovered that resonance can be
learned from co-occurrence. Vaswani and collaborators discovered that
resonance can be learned non-symmetrically and applied dynamically.
Each discovery, taken in isolation, looked like a refutation of its
predecessors. Taken together, they describe the operators and the
sub-algebras of a single idea algebra. The next chapter takes up the
question of how this algebra is realized in the brain, where the same
mathematical structure must be implemented in spike trains, synaptic
weights, and population codes.
